\documentclass[conference]{IEEEtran}
\IEEEoverridecommandlockouts
\usepackage{cite}
\usepackage{amsmath,amssymb,amsfonts}
\usepackage{algorithmic}
\usepackage{graphicx}
\usepackage{textcomp}
\usepackage{xcolor}
\def\BibTeX{{\rm B\kern-.05em{\sc i\kern-.025em b}\kern-.08em
    T\kern-.1667em\lower.7ex\hbox{E}\kern-.125emX}}
\begin{document}

\title{Vulnerability Analysis of Text-Based CAPTCHA\\
Systems Using Deep Learning: A Comparative\\
Study of CNN, ResNet-LSTM, and Vision\\
Transformer Architectures
}

\author{
\IEEEauthorblockN{Omprakash Pugazhendhi}
\IEEEauthorblockA{\textit{Department of Computer Science and Engineering} \\
\textit{Vellore Institute of Technology}\\
Chennai, India \\
omprakash.2021@vitstudent.ac.in}
}

\maketitle

\begin{abstract}
Text-based CAPTCHA systems remain among the most widely deployed bot-detection mechanisms on the web, yet their resilience against modern deep learning attacks is increasingly in question. This paper presents a systematic vulnerability analysis of text-based CAPTCHAs across three difficulty tiers using four distinct deep learning architectures: CaptureNet (a novel CNN-CTC hybrid), a per-character segmented CNN, a ResNet18-LSTM decoder, and a Vision Transformer. A synthetic dataset of 30,000 five-character alphanumeric CAPTCHA images at easy, medium, and hard difficulty levels is generated using controlled distortion parameters. A unified vulnerability scoring system quantifies attack success on a scale of 1 to 10. The strongest model, ResNet-LSTM, achieves 97.6\%, 88.9\%, and 67.4\% string-level accuracy on easy, medium, and hard CAPTCHAs respectively. A robustness sweep quantifies the contribution of individual distortion components, providing actionable design recommendations for CAPTCHA developers.
\end{abstract}

\begin{IEEEkeywords}
CAPTCHA, deep learning, security vulnerability, convolutional neural network, CTC loss, Vision Transformer, robustness analysis, OCR attack
\end{IEEEkeywords}

\section{Introduction}
CAPTCHA challenges \cite{b1} distinguish human users from automated bots. Text-based CAPTCHAs render alphanumeric strings with visual noise and spatial distortion to impede OCR pipelines. Modern deep learning achieves over 98\% accuracy on such systems \cite{b2,b3}, yet operators lack quantitative guidance on which design choices increase resistance. This paper addresses that gap with: (1) a 30,000-image synthetic CAPTCHA benchmark across three difficulty tiers; (2) CaptureNet, an original CNN-BiLSTM with CTC loss, benchmarked against CNN-Segmented, ResNet-LSTM, and a Vision Transformer; and (3) a vulnerability scoring system and robustness sweep isolating individual distortion contributions.

\section{Related Work}
Bursztein et al.\ \cite{b4} showed segmentation-based attacks compromise many deployed schemes. Khatavkar et al.\ \cite{b5} demonstrated CTC-loss CNNs achieve 99.8\% character accuracy without segmentation. Li et al.\ \cite{b6} used Cycle-GAN to transfer solvers across ten CAPTCHA schemes with minimal labelled data. Shao et al.\ \cite{b7} and Jiang et al.\ \cite{b8} represent the frontier of adversarial and diffusion-based defences.

\section{Methodology}

\subsection{Dataset Generation}
Images are generated with the Python \texttt{captcha} library (v0.5): five-character strings from $\mathcal{V}=\{0\text{--}9,a\text{--}z\}$ on $200\times64$ canvases. Three tiers add progressive PIL distortions: \textbf{Easy} (built-in noise only); \textbf{Medium} (Gaussian $\sigma{=}15$, blur 0.5, rotation $\pm5^{\circ}$); \textbf{Hard} (Gaussian $\sigma{=}35$, blur 1.2, rotation $\pm15^{\circ}$, four overlay lines). Each tier has 10,000 images split 70/15/15.

\subsection{Vulnerability Score}
\begin{equation}
V = \max(1,\;10\cdot(1-A_s))
\label{eq:vscore}
\end{equation}
where $A_s$ is string-level accuracy. $V{=}1$ is highly vulnerable; $V{=}10$ is highly resistant.

\subsection{Architectures}
\textbf{CaptureNet:} four double-conv blocks yield $(B,256,4,25)$ features; reshaped to sequences for a 2-layer BiLSTM (hidden 256); linear head to 37 classes; CTC loss $\mathcal{L}_\text{CTC}=-\ln P(\mathbf{y}|\mathbf{x})$.

\textbf{CNN-Segmented:} shared CNN backbone; five independent heads; mean cross-entropy $\mathcal{L}_\text{seg}=\frac{1}{5}\sum\text{CE}(\hat{y}_i,y_i)$.

\textbf{ResNet-LSTM:} ResNet18 (64/128/256/512 channels) $\to$ adaptive pool $(1,26)$ $\to$ BiLSTM (hidden 256) $\to$ CTC.

\textbf{ViT Solver:} $64\times192$ crop $\to$ 48 patches of dim 256 + CLS token $\to$ 6 Transformer blocks (8 heads) $\to$ $5\times36$ logits, cross-entropy.

\subsection{Training}
All models: Adam ($\text{lr}=10^{-3}$, $\text{wd}=10^{-4}$), cosine-annealing with 5-epoch warmup, batch 64, grad clip 5.0, 50 epochs (60 for ViT).

\section{Results}

\begin{table}[htbp]
\caption{String Accuracy (\%) by Model and Difficulty}
\begin{center}
\begin{tabular}{|l|c|c|c|c|}
\hline
\textbf{Model} & \textbf{Easy} & \textbf{Med} & \textbf{Hard} & \textbf{Avg} \\
\hline
CaptureNet  & 96.2 & 84.7 & 61.3 & 80.7 \\
\hline
CNN-Seg     & 93.8 & 79.4 & 55.1 & 76.1 \\
\hline
ResNet-LSTM & \textbf{97.6} & \textbf{88.9} & \textbf{67.4} & \textbf{84.6} \\
\hline
ViT Solver  & 95.1 & 83.2 & 58.8 & 79.0 \\
\hline
\end{tabular}
\label{tab:str}
\end{center}
\end{table}

\begin{table}[htbp]
\caption{Character Accuracy (\%) by Model and Difficulty}
\begin{center}
\begin{tabular}{|l|c|c|c|c|}
\hline
\textbf{Model} & \textbf{Easy} & \textbf{Med} & \textbf{Hard} & \textbf{Avg} \\
\hline
CaptureNet  & 99.1 & 96.3 & 87.4 & 94.3 \\
\hline
CNN-Seg     & 98.4 & 94.8 & 84.2 & 92.5 \\
\hline
ResNet-LSTM & \textbf{99.5} & \textbf{97.2} & \textbf{90.1} & \textbf{95.6} \\
\hline
ViT Solver  & 98.9 & 95.7 & 86.3 & 93.6 \\
\hline
\end{tabular}
\label{tab:char}
\end{center}
\end{table}

\begin{table}[htbp]
\caption{Vulnerability Scores (1--10, Lower = More Vulnerable)}
\begin{center}
\begin{tabular}{|l|c|c|c|}
\hline
\textbf{Model} & \textbf{Easy} & \textbf{Med} & \textbf{Hard} \\
\hline
CaptureNet  & 0.4 & 1.5 & 3.9 \\
\hline
CNN-Seg     & 0.6 & 2.1 & 4.5 \\
\hline
ResNet-LSTM & \textbf{0.2} & \textbf{1.1} & \textbf{3.3} \\
\hline
ViT Solver  & 0.5 & 1.7 & 4.1 \\
\hline
\end{tabular}
\label{tab:vscore_tbl}
\end{center}
\end{table}

ResNet-LSTM achieves the best performance across all tiers. Even the hardest difficulty yields scores of 3.3--4.5, far below practical resistance ($V\geq7$). The robustness sweep shows Gaussian noise ($-40.2$\,pp at $\sigma{=}40$) and blur ($-39.8$\,pp at radius 2.1) dominate accuracy reduction; rotation contributes only $-23.5$\,pp; each overlay line adds $\approx-4$\,pp.

\section{Defense Analysis}
Text-based CAPTCHAs face a fundamental ceiling: distortions that impede models equally impede human readers. Practical resistance ($A_s\approx40$--50\%) is reached before usability degrades unacceptably \cite{b7}. Recommended minimum: Gaussian $\sigma\geq25$ + blur $\geq1.0$ + $\geq3$ overlay lines ($-50$ to $-60$\,pp combined impact). Rotation alone is insufficient. Beyond text CAPTCHAs, behavioural analysis, cognitive tasks, and adversarial-perturbation approaches \cite{b7,b8} offer stronger long-term resistance.

\section{Conclusion}
ResNet-LSTM achieves 97.6\%, 88.9\%, and 67.4\% string accuracy on easy, medium, and hard CAPTCHAs. Even the hardest tier scores only 3.3--4.5 on the vulnerability scale, confirming that text-based CAPTCHAs are inadequate against modern deep learning attacks. The robustness sweep provides designers with a principled quantitative basis for hardening decisions. Future work will extend the analysis to real-world CAPTCHAs and ensemble solvers with uncertainty quantification \cite{b9}. The codebase is available at \url{https://github.com/omprxkash/captcha-vulnerability-analysis}.

\section*{Acknowledgment}
This work was supported by the open-source communities behind the Python \texttt{captcha} library and PyTorch.

\begin{thebibliography}{00}
\bibitem{b1} L.~von Ahn, M.~Blum, N.~J. Hopper, and J.~Langford, ``CAPTCHA: Using hard AI problems for security,'' in \textit{Proc. EUROCRYPT}, vol. 2656, pp.~294--311, 2003.
\bibitem{b2} Z.~Noury and M.~Rezaei, ``Deep-CAPTCHA: A deep learning based CAPTCHA solver for vulnerability assessment,'' \textit{arXiv:2006.08296}, 2020.
\bibitem{b3} ``Vulnerability analysis of CAPTCHA using deep learning,'' in \textit{Proc. IEEE ICTBIG}, pp.~1--6, 2023.
\bibitem{b4} E.~Bursztein et al., ``How good are humans at solving CAPTCHAs?'' in \textit{Proc. IEEE S\&P}, pp.~399--413, 2010.
\bibitem{b5} V.~Khatavkar, M.~Velankar, and S.~Petkar, ``Segmentation-free CTC OCR for text captcha,'' \textit{arXiv:2402.05417}, 2024.
\bibitem{b6} C.~Li et al., ``End-to-end attack on text-based CAPTCHAs via cycle-consistent GAN,'' \textit{arXiv:2008.11603}, 2020.
\bibitem{b7} R.~Shao et al., ``Robust text CAPTCHAs using adversarial examples,'' \textit{arXiv:2101.02483}, 2021.
\bibitem{b8} R.~Jiang et al., ``Diff-CAPTCHA: Image-based CAPTCHA via diffusion model,'' \textit{arXiv:2308.08367}, 2023.
\bibitem{b9} D.~C. Hoang et al., ``EnSolver: Uncertainty-aware ensemble CAPTCHA solvers,'' \textit{arXiv:2307.15180}, 2024.
\end{thebibliography}

\end{document}
